\chapter{Claim Development and LDFs}
\label{chap:claim_development}

\section{Overview}

Claims do not typically reach their ultimate cost on the date they are
reported. Reserving actuaries study \emph{development patterns} that
describe how reported and paid losses evolve from the date of an
accident through to ultimate settlement. \actsim provides a simple,
transparent mechanism for developing simulated claims using user-supplied
loss development factors (LDFs).

\section{Definitions}

The terminology used in this chapter follows standard actuarial
practice.

\begin{description}[leftmargin=*,style=nextline]
    \item[Accident year] The calendar year in which a claim occurred.
    \item[Development age] The number of months elapsed since the start
        of the accident period (often measured at quarterly or annual
        intervals).
    \item[Reported loss] The sum of paid losses and case reserves on a
        claim at a given valuation date.
    \item[Paid loss] The cumulative amount paid on a claim at a given
        valuation date.
    \item[Incurred loss] In this manual, used interchangeably with
        reported loss.
    \item[Ultimate loss] The expected total cost of a claim once all
        development is complete.
    \item[Loss development factor (LDF)] The ratio of losses at one
        development age to losses at the previous age.
    \item[Cumulative development factor (CDF)] The ratio of ultimate
        losses to losses at a given development age.
\end{description}

\section{Mathematical Setup}

Let $U$ denote the ultimate loss for a claim and let $\mathrm{LDF}_a$
denote the development factor between age $a$ and the next reported age.
The cumulative development factor at age $a$ is
\[
    \mathrm{CDF}_a = \prod_{j \geq a} \mathrm{LDF}_j,
\]
where the product runs over all development ages from $a$ through to
ultimate. The reported loss at age $a$ is then
\[
    L_a = \frac{U}{\mathrm{CDF}_a}.
\]
The incremental development between two adjacent ages is
\[
    \Delta L_a = L_a - L_{a-1}.
\]

\section{Stochastic LDFs in ActSim}

In the deterministic chain-ladder method, the LDF at each development
age is a single point estimate. \actsim introduces stochastic variation
around the user-supplied base LDFs by multiplying each factor by a
normal random variable centred at $1$:
\[
    \widetilde{\mathrm{LDF}}_a = \mathrm{LDF}_a \cdot \varepsilon_a,
    \qquad \varepsilon_a \sim \mathcal{N}(1, \sigma^2).
\]
The standard deviation $\sigma$ is the \texttt{volatility} argument
passed to \texttt{simulate\_claim\_development}. A separate set of
stochastic LDFs is drawn for each claim, so different claims at the
same accident age develop along different paths.

\section{Tail Factor}

The \texttt{cumulative\_factor} argument represents an additional
factor applied beyond the user-supplied LDF schedule, intended to
capture residual development beyond the latest age. For a fully
developed pattern, set the tail factor to $1.0$.

\section{Worked Calculation}

Suppose a claim has ultimate loss $U = 10{,}000$ and the user has
specified the following deterministic LDF schedule.

\begin{center}
\begin{tabular}{rrr}
\toprule
\textbf{Age (months)} & \textbf{LDF} & \textbf{CDF} \\
\midrule
 0  & 2.00 & 4.987 \\
 3  & 1.50 & 2.494 \\
 6  & 1.20 & 1.663 \\
 9  & 1.10 & 1.386 \\
12  & 1.05 & 1.260 \\
15  & 1.02 & 1.200 \\
18  & 1.00 & 1.000 \\
\bottomrule
\end{tabular}
\end{center}

The reported loss at each age is then $L_a = U / \mathrm{CDF}_a$, giving
the development pattern shown below (rounded).

\begin{center}
\begin{tabular}{rr}
\toprule
\textbf{Age (months)} & \textbf{Reported loss} \\
\midrule
 0  & 2{,}005 \\
 3  & 4{,}010 \\
 6  & 6{,}013 \\
 9  & 7{,}215 \\
12  & 7{,}937 \\
15  & 8{,}333 \\
18  & 10{,}000 \\
\bottomrule
\end{tabular}
\end{center}

In a stochastic simulation, the reported losses at each age will vary
from claim to claim, but their expected pattern matches the values
above.

\section{From Claim Development to Triangles}

The long-format \texttt{claim\_development} DataFrame produced by
\actsim contains one row per (claim, development age). To form a
loss-development triangle, the DataFrame must be aggregated by
accident period and development age. This is straightforward in pandas:

\begin{lstlisting}[style=pythonstyle]
triangle = (
    claim_sim.claim_development
    .groupby(["accident_year", "development_month"])["incurred_loss"]
    .sum()
    .reset_index()
)
\end{lstlisting}

This long-format triangle is the canonical input to the
\texttt{chainladder} package (see Chapter~\ref{chap:chainladder}).

\section{Practical Considerations}

\begin{itemize}[leftmargin=*]
    \item Choose an LDF schedule that is consistent with the simulated
          claim mix and the calendar period of the simulation. Using
          patterns that are inconsistent with the underlying claim
          process will produce triangles that disagree with theoretical
          ultimates.
    \item The \texttt{volatility} parameter controls the dispersion of
          development paths but not the central tendency. For
          well-defined ultimates, the simulated triangle should
          converge to the deterministic pattern as the number of claims
          increases.
    \item For lines of business with long tails, the user should
          extend the LDF schedule to cover the full development horizon
          and consider whether a tail factor is required.
\end{itemize}
